\documentclass[11pt,a4paper]{ctexart}
\usepackage[a4paper,margin=2.2cm]{geometry}
\usepackage{amsmath,amssymb}
\usepackage{booktabs}
\usepackage{longtable}
\usepackage{array}
\usepackage{enumitem}
\usepackage{hyperref}
\usepackage{xcolor}
\usepackage{listings}

\hypersetup{
  colorlinks=true,
  linkcolor=blue!60!black,
  urlcolor=blue!60!black,
  pdftitle={CGYRO Comparison 模块完整说明书},
  pdfauthor={CGYRO Team}
}

\setlist[itemize]{leftmargin=1.8em}
\setlist[enumerate]{leftmargin=2.0em}
\setlist[description]{leftmargin=2.2em}

\lstset{
  basicstyle=\ttfamily\small,
  breaklines=true,
  frame=single,
  columns=fullflexible
}

\title{CGYRO Comparison 模块完整说明书}
\author{代码路径：\texttt{CGYRO/cgyro\_comparison*.py}}
\date{\today}

\begin{document}
\maketitle
\tableofcontents
\newpage

\section{文档说明}
\subsection{目的}
本文档用于完整说明 \texttt{cgyro\_comparison} 模块的：
\begin{itemize}
  \item 架构设计与代码组织；
  \item GUI 功能与每个选项的行为；
  \item 绘图后端算法、输入数据要求与输出；
  \item 配置项、路径解析、容错与常见报错排查；
  \item 二次开发扩展流程与建议。
\end{itemize}

\subsection{适用读者}
\begin{itemize}
  \item 使用 GUI 进行多案例对比绘图的使用者；
  \item 维护/扩展模块的开发者；
  \item 需要把结果用于报告或后续自动化流程的用户。
\end{itemize}

\subsection{覆盖文件}
\begin{itemize}
  \item \texttt{cgyro\_comparison.py}
  \item \texttt{cgyro\_comparison\_bootstrap.py}
  \item \texttt{cgyro\_comparison\_ui.py}
  \item \texttt{cgyro\_comparison\_plotting.py}
\end{itemize}

\section{模块总览}
\subsection{总体架构}
主类为：
\begin{center}
\texttt{CGYRO\_Comparison(CgyroUiMixin, CgyroPlottingMixin)}
\end{center}

\begin{description}
  \item[\texttt{CgyroUiMixin}] 负责窗口与控件构建、案例管理、参数联动、用户交互。
  \item[\texttt{CgyroPlottingMixin}] 负责数据处理、算法实现、图形渲染与后端分发。
\end{description}

\subsection{入口流程}
\texttt{cgyro\_comparison.py} 中的运行流程：
\begin{enumerate}
  \item 创建 Tk 根窗口 \texttt{tk.Tk()}；
  \item 初始化 \texttt{CGYRO\_Comparison(root)}；
  \item 进入事件循环 \texttt{root.mainloop()}。
\end{enumerate}

\subsection{运行时对象}
\begin{itemize}
  \item \texttt{self.cases}: \texttt{dict[str, data\_obj]}，已加载案例缓存。
  \item \texttt{self.fig / self.ax}: 当前 Matplotlib 图对象。
  \item \texttt{self.ani}: 动画句柄，控制播放/暂停/逐帧。
\end{itemize}

\section{依赖与环境}
\subsection{必需依赖}
\begin{itemize}
  \item Python 3
  \item \texttt{numpy}
  \item \texttt{matplotlib}
  \item \texttt{tkinter}
  \item \texttt{pygacode}（核心数据读取）
\end{itemize}

\subsection{可选依赖}
\begin{itemize}
  \item \texttt{scipy.signal} 与 \texttt{scipy.optimize}：用于 \texttt{rcorr\_phi} 的 Hilbert 包络和指数拟合。
\end{itemize}
未安装 SciPy 时，相关拟合步骤会自动降级，不影响其他图型。

\subsection{pygacode 导入引导}
\texttt{cgyro\_comparison\_bootstrap.py} 会尝试向 \texttt{sys.path} 注入：
\begin{itemize}
  \item \texttt{../gacode/f2py}
  \item \texttt{../gacode-master/f2py}
\end{itemize}
然后导入：
\begin{itemize}
  \item \texttt{pygacode.cgyro.data.cgyrodata}
  \item \texttt{pygacode.cgyro.data\_plot.cgyrodata\_plot}
\end{itemize}

\subsection{fallback mock}
若 pygacode 不可用，会创建最小 mock 类用于 GUI 调试。注意：
\begin{itemize}
  \item mock 仅用于界面联调；
  \item 不代表真实物理结果。
\end{itemize}

\section{默认配置与常量}
\begin{longtable}{>{\ttfamily}p{0.35\linewidth} p{0.57\linewidth}}
\toprule
常量名 & 说明 \\
\midrule
\endhead
DEFAULT\_APP\_TITLE & 窗口标题，默认 \texttt{"CGYRO Comparison Tool"}。 \\
DEFAULT\_WINDOW\_GEOMETRY & 窗口初始大小，默认 \texttt{"1200x800"}。 \\
DEFAULT\_CASE\_PICKER\_ROOT & 案例选择初始目录，默认 \texttt{/data/share/<当前用户名>}。 \\
DEFAULT\_LINEAR\_GAMMA\_FILE & 线性参考文件默认名 \texttt{omega\_gamma\_vs\_ky.txt}。 \\
DEFAULT\_EXPORT\_DIRNAME & 默认导出目录名 \texttt{CGYRO\_vs\_CGYRO\_exports}。 \\
\bottomrule
\end{longtable}

\section{安装与启动}
\subsection{推荐启动方式}
在 \texttt{CGYRO} 目录执行：
\begin{lstlisting}
python cgyro_comparison.py
\end{lstlisting}

\subsection{模块导入方式}
代码支持两种导入风格：
\begin{itemize}
  \item 脚本式：\texttt{from cgyro\_comparison\_ui import ...}
  \item 包式：\texttt{from .cgyro\_comparison\_ui import ...}
\end{itemize}

\section{GUI 详细说明}
\subsection{布局}
\begin{itemize}
  \item 左侧：案例管理 + 绘图控制 + 动画按钮。
  \item 右侧：Matplotlib 画布 + toolbar。
\end{itemize}

\subsection{File 菜单}
\begin{itemize}
  \item \texttt{Add Case (Single)}
  \item \texttt{Add Case (Multiple)}
  \item \texttt{Add Group}
  \item \texttt{Exit}
\end{itemize}

\subsection{案例管理区}
\begin{description}
  \item[\texttt{Add Case (Single)}] 直接选择一个案例目录并加载，适合逐个补充。
  \item[\texttt{Add Case (Multiple)}] 先选父目录，再在子目录列表中 Ctrl/Shift 多选，批量加载多个案例（推荐用于扫描对比）。
  \item[\texttt{Add Group}] 自动加载父目录下所有“疑似案例目录”，适合全量导入。
  \item[\texttt{Remove}] 删除当前选中案例。
  \item[\texttt{Remove All}] 清空案例缓存并确认。
  \item[\texttt{Reload}] 使用原路径重载数据对象。
\end{description}

\subsection{Add Case (Multiple) 推荐用法（重点）}
\begin{enumerate}
  \item 点击 \texttt{Add Case (Multiple)}。
  \item 先选择“父目录”（该目录下应包含多个 case 子目录）。
  \item 在弹窗列表中使用 Ctrl/Shift 选择需要比较的多个子目录。
  \item 点击确认后一次性完成批量加载。
  \item 在左侧案例列表中拖拽排序，再执行 \texttt{Plot} 进行对比。
\end{enumerate}
\noindent\textbf{建议：} 当你做参数扫描、不同物理模型对比、同一 shot 多阶段对比时，优先使用 \texttt{Add Case (Multiple)}，效率明显高于逐个单独添加。

\subsection{案例目录判定规则}
若目录包含以下任意文件，则视作候选案例：
\begin{itemize}
  \item \texttt{input.cgyro}
  \item \texttt{out.cgyro.freq}
  \item \texttt{input.gacode}
\end{itemize}

\subsection{案例命名冲突}
同名目录会自动改名为 \texttt{name\_1, name\_2, ...}，保证 \texttt{self.cases} 键唯一。

\subsection{案例列表拖拽重排}
支持鼠标拖拽改变案例顺序。该顺序也会影响默认绘图顺序和图例顺序。

\section{Plot Type 与分发机制}
\subsection{GUI 到内部 token 的转换}
\texttt{\_build\_effective\_plot\_type()} 将 GUI 选项映射为内部字符串。

\begin{longtable}{p{0.22\linewidth} p{0.28\linewidth} p{0.42\linewidth}}
\toprule
GUI 主类型 & 内部 token 示例 & 备注 \\
\midrule
\endhead
Frequency & \texttt{Frequency} & 支持 \texttt{Divided by ky} \\
Growth Rate & \texttt{Growth Rate} & 支持 \texttt{Divided by ky} \\
Flux & \texttt{Energy Flux vs ky} 等 & 支持 Decomp、Species \\
Fluctuation 1D & \texttt{Phi vs ky}, \texttt{Apar FFT} & FFT 为特殊分支 \\
Fluctuation 2D & \texttt{Fluctuation 2D} / \texttt{Fluctuation 2D vs xt} & 轮廓类图 \\
Zonal ExB Shearing Rate & \texttt{ZF ExB Shearing Rate} 等 & 三种子模式 \\
Others & \texttt{Integration Error} 等 & Error/rcorr/POD \\
\bottomrule
\end{longtable}

\subsection{单案例限制}
以下图型仅支持单案例：
\begin{itemize}
  \item FFT（\texttt{* FFT}）
  \item \texttt{Fluctuation 2D} / \texttt{Fluctuation 2D vs xt}
  \item \texttt{POD Parity}
\end{itemize}
若选中多个案例，程序会告警并只绘制第一个。

\subsection{后端分发}
\texttt{CgyroPlottingMixin.\_dispatch\_single\_case\_plot()} 按 token 路由到具体函数：
\begin{itemize}
  \item \texttt{\_plot\_frequency\_growth}
  \item \texttt{\_plot\_flux}
  \item \texttt{\_plot\_fluctuation\_1d}
  \item \texttt{\_plot\_fluctuation\_fft}
  \item \texttt{\_plot\_fluctuation\_2d / \_plot\_xt\_fluctuation\_contours}
  \item \texttt{\_plot\_zf\_exb\_shearing\_*}
  \item \texttt{\_plot\_other\_error / \_plot\_other\_rcorr\_phi / \_plot\_other\_apar\_pod\_parity}
\end{itemize}

\section{动态选项全量说明}
\subsection{通用选项}
\begin{description}
  \item[Time Start / Time End] 用于时间平均区间、统计区间、动画帧区间。
  \item[Clear Time] 清空时窗，回退自动策略（后 50\%）。
  \item[Log X / Log Y] 最终线图坐标轴缩放。
\end{description}

\subsection{Flux 选项}
\begin{itemize}
  \item \texttt{Flux Type}: Energy / Particle。
  \item \texttt{X-axis}: \texttt{vs ky}, \texttt{vs kx (estimated)}, \texttt{vs Time}。
  \item \texttt{Species}: 单一物种 / Main Ion(D+T) / All Ions。
  \item \texttt{Decompose by Field}: 分解 ES/EM 通道。
  \item \texttt{Plot All Species (First Case Only)}: 第一案例多物种叠加。
\end{itemize}

\subsection{Fluctuation 1D 选项}
\begin{itemize}
  \item 字段：\texttt{Phi}, \texttt{Apar}, \texttt{Bpar}
  \item X 轴：\texttt{vs ky}, \texttt{vs kx}, \texttt{vs Time}, \texttt{fft}
\end{itemize}

\subsection{FFT 选项}
\begin{description}
  \item[Mode] \texttt{Linear}（coherent）或 \texttt{Nonlinear}（incoherent）。
  \item[View] \texttt{Omega vs ky} 或 \texttt{Omega vs kx}。
  \item[Spectrum] \texttt{Amplitude} 或 \texttt{Power}。
  \item[Overlay Linear Freq] 叠加线性参考频率。
  \item[Linear File] 线性参考文件路径。
\end{description}

\subsection{Fluctuation 2D 选项}
\begin{itemize}
  \item View：\texttt{vs xy}（动画）、\texttt{vs xt}（时空轮廓）。
  \item Moment：\texttt{Phi, Density, Energy, Temperature, Apar, Bpar}。
  \item 当 moment 为 \texttt{Density/Energy/Temperature}，需选择 Species。
  \item 勾选 \texttt{Spatial axes normalize to electron scale (x,y)/\rho_e} 后：
  空间坐标按电子尺度显示。对于 \texttt{vs xy}，x/y 同时从 $\rho_s$ 转为 $\rho_e$；
  对于 \texttt{vs xt}，纵轴 $x$ 从 $\rho_s$ 转为 $\rho_e$。
\end{itemize}

\subsection{Zonal ExB 选项}
\begin{itemize}
  \item \texttt{vs Time}
  \item \texttt{vs kx}
  \item \texttt{vs Gamma\_Linear}（需要线性文件）
\end{itemize}

\subsection{Others 选项}
\begin{itemize}
  \item \texttt{Error}
  \item \texttt{rcorr\_phi}: 字段（Phi/Apar/Bpar） + Theta 索引
  \item \texttt{POD\_parity}: 字段（Apar/Phi） + kx/ky 选择
\end{itemize}

\section{物种选择逻辑}
\subsection{候选集合}
\texttt{\_update\_species\_list()} 会计算“所有已加载案例的公共物种集合”，并输出到下拉框。

\subsection{特殊选项}
\begin{itemize}
  \item \texttt{Main Ion (D+T)}：按物种名识别 Deuterium/Tritium。
  \item \texttt{All Ions}：按电荷 \(Z>0\) 识别所有离子。
\end{itemize}

\subsection{物种名映射}
根据 \((Z,M)\) 近似映射为 Electron/Hydrogen/Deuterium/...；无法识别时显示
\texttt{Ion (Z=?, M=?)}。

\section{时间窗口与自动回退}
\texttt{\_get\_time\_indices(t)} 行为：
\begin{enumerate}
  \item 空输入：默认后 50\% 时间段。
  \item 超范围输入：钳制到 \([t_{\min},t_{\max}]\)。
  \item 起止反向或重合：回退后 50\%。
  \item 仍无有效索引：使用后半段索引作为兜底。
\end{enumerate}

\section{主要绘图后端详解}
\subsection{Frequency / Growth Rate}
\begin{itemize}
  \item 数据源：\texttt{data.freq}, \texttt{data.ky}
  \item 支持二维/三维布局的 \texttt{freq} 自动识别；
  \item 时间窗有效时做均值；
  \item 可选除以 \(k_y\)。
\end{itemize}

\subsection{Flux}
\subsubsection{vs ky / vs Time}
\begin{itemize}
  \item 数据源：\texttt{data.ky\_flux}
  \item 若未加载，自动尝试 \texttt{data.getflux()}。
  \item 维度预期：5D（species, moment, field, ky, t）。
\end{itemize}

\subsubsection{vs kx (estimated)}
\begin{itemize}
  \item 不依赖 \texttt{ky\_flux}；
  \item 从场与矩估算 ES/EM 分量并求和：
  \[
  \Gamma_{tot}=\Gamma_{ES}+\Gamma_{EM,A}+\Gamma_{EM,B}.
  \]
  \item 仅绘制 \(k_x \ge 0\) 支路。
\end{itemize}

\subsection{Fluctuation 1D}
通过 \texttt{kxky\_phi/apar/bpar}（必要时触发 \texttt{getbigfield()}）构建：
\begin{itemize}
  \item vs ky：对 \(k_x\) 求和并在时窗内平均；
  \item vs kx：对 \(k_y\) 求和并在时窗内平均；
  \item vs Time：输出 \(n=0\) 与 \(n>0\) 两条 RMS 时间曲线。
\end{itemize}

\subsection{Fluctuation FFT}
\begin{itemize}
  \item 支持 \( \omega\)-\(k_y \) 与 \( \omega\)-\(k_x \) 视图；
  \item 支持 Amplitude / Power；
  \item Linear 模式使用 coherent 叠加，Nonlinear 模式使用 incoherent 叠加；
  \item 可叠加线性参考频率曲线。
\end{itemize}

\subsection{Fluctuation 2D}
\begin{description}
  \item[\texttt{vs xy}] 根据时间索引生成动画，支持暂停/单步。
  \item[\texttt{vs xt}] 先重构 \(x\) 方向剖面，再画时空轮廓。
\end{description}

\subsection{Zonal ExB}
核心定义：
\[
\Omega(k_x,t)=-k_x^2\phi_{ZF}(k_x,t),\qquad
\omega_{ZF}(t)=\sqrt{\sum_{k_x}|\Omega(k_x,t)|^2}.
\]
三种输出：
\begin{itemize}
  \item vs Time：\(\omega_{ZF}(t)\)；
  \item vs kx：\(\Omega(k_x,t)\) 的时窗均值谱；
  \item vs Gamma\_Linear：\(\langle\omega_{ZF}\rangle_t\) 与 \(\gamma_{lin}(k_y)\) 同图。
\end{itemize}

\subsection{Others: Error / rcorr\_phi / POD parity}
\begin{itemize}
  \item Error：绘制 \texttt{err1/err2}，并默认 y 轴对数。
  \item rcorr\_phi：径向相关函数；若 SciPy 可用则做包络拟合，输出相关长度。
  \item POD parity：SVD 模态奇偶性判据
  \[
  P_n=\frac{|\int a_n(z)\,dz|}{\int |a_n(z)|\,dz}.
  \]
  分类阈值：\(P>0.7\) tearing-like，\(P<0.3\) ballooning-like，其余 mixed。
\end{itemize}

\section{线性参考文件解析}
\subsection{文件格式}
每一行取前 3 个浮点数作为：
\[
k_y,\ \omega,\ \gamma.
\]
注释行以 \texttt{\#} 开头。

\subsection{路径查找优先级}
\texttt{\_resolve\_linear\_gamma\_file\_path()} 查找顺序：
\begin{enumerate}
  \item 若是绝对路径：直接检查；
  \item 案例目录；
  \item 案例目录上一级；
  \item 当前工作目录；
  \item \texttt{\$CWD/DEFAULT\_EXPORT\_DIRNAME/<filename>}。
\end{enumerate}

\section{数据形状与兼容策略}
\subsection{核心目标}
不同版本/来源的数据形状可能不一致，模块通过 \texttt{\_coerce\_kxky\_complex()} 和
\texttt{\_load\_kxky\_complex()} 做统一处理。

\subsection{支持的典型形状}
\begin{itemize}
  \item 复数 4D：\([nr,\theta,ky,t]\)
  \item 复数 5D：\([nr,\theta,species,ky,t]\)
  \item 打包实虚 5D：\([2,nr,\theta,ky,t]\)
  \item 打包实虚 6D：\([2,nr,\theta,species,ky,t]\)
\end{itemize}

\subsection{失败处理}
若形状无法识别：
\begin{itemize}
  \item 打印 \texttt{Unsupported ... shape}；
  \item 跳过当前案例该图型；
  \item GUI 保持可继续使用。
\end{itemize}

\section{容错与提示体系}
\subsection{提示方式}
\begin{itemize}
  \item 终端 \texttt{print}: 数据缺失、形状不符、拟合失败、回退行为等。
  \item \texttt{messagebox.showwarning/showerror}: 影响用户决策的提示（如单案例限制、theta 分辨率不足）。
\end{itemize}

\subsection{典型容错行为}
\begin{itemize}
  \item 大字段缺失时自动尝试 \texttt{getbigfield()}。
  \item \texttt{ky\_flux} 缺失时自动尝试 \texttt{getflux()}。
  \item 时间窗口非法时自动回退。
  \item 线性文件缺失时给出清晰提示并返回，不中断 GUI。
\end{itemize}

\section{性能与使用建议}
\begin{itemize}
  \item 大案例 + 2D 动画 + FFT 同时操作时内存压力较大，建议分步骤操作。
  \item 多案例对比优先使用线图（Frequency/Growth/Flux），再对单案例做 2D/FFT 深入分析。
  \item 需要重复分析时，优先使用 \texttt{Reload} 而不是重新选择目录。
\end{itemize}

\section{开发者扩展指南}
\subsection{新增图型标准流程}
\begin{enumerate}
  \item 在 \texttt{CgyroUiMixin.\_init\_options()} 创建控件；
  \item 在 \texttt{update\_options()} 控制显隐；
  \item 在 \texttt{\_build\_effective\_plot\_type()} 添加 token 映射；
  \item 在 \texttt{\_dispatch\_single\_case\_plot()} 注册 handler；
  \item 在 \texttt{\_apply\_axis\_labels()} 补默认标签；
  \item 必要时增加数据加载/形状兼容函数。
\end{enumerate}

\subsection{推荐代码原则}
\begin{itemize}
  \item 函数级职责单一，复用已有 helper；
  \item 对数组维度、有限值、空数据做显式检查；
  \item 遵循“可回退、不中断 GUI”原则；
  \item 失败信息要可定位（案例名 + 函数语义）。
\end{itemize}

\section{函数索引（核心）}
\subsection{UI 侧}
\begin{longtable}{>{\ttfamily}p{0.33\linewidth} p{0.58\linewidth}}
\toprule
函数 & 作用 \\
\midrule
\endhead
\_create\_layout & 创建主窗口、左侧控件、右侧画布与 toolbar。 \\
\_init\_options & 初始化所有动态选项控件。 \\
update\_options & 根据 Plot Type 刷新选项布局。 \\
add\_case / add\_group & 加载单个或批量案例目录。 \\
reload\_cases & 按原路径重载案例对象。 \\
\_build\_effective\_plot\_type & GUI 选项到内部 token 的转换。 \\
\_resolve\_species\_indices & 将物种文本选择解析为索引列表。 \\
\bottomrule
\end{longtable}

\subsection{Plotting 侧}
\begin{longtable}{>{\ttfamily}p{0.33\linewidth} p{0.58\linewidth}}
\toprule
函数 & 作用 \\
\midrule
\endhead
plot\_comparison & 主绘图入口，处理单案例限制和最终渲染。 \\
\_dispatch\_single\_case\_plot & 内部 token 到后端函数的路由。 \\
\_plot\_frequency\_growth & 频率/增长率绘图。 \\
\_plot\_flux & Flux 三模式绘图主入口。 \\
\_plot\_flux\_vs\_kx\_estimated & 估算型 Flux-kx 谱。 \\
\_plot\_fluctuation\_1d & 一维涨落图。 \\
\_plot\_fluctuation\_fft & 频谱图。 \\
\_plot\_fluctuation\_2d & xy 动画。 \\
\_plot\_xt\_fluctuation\_contours & xt 轮廓图。 \\
\_plot\_zf\_exb\_shearing\_* & Zonal ExB 三子模式。 \\
\_plot\_other\_error / rcorr / pod & Others 三子模式。 \\
\_coerce\_kxky\_complex & 统一复杂数组形状。 \\
\_load\_kxky\_complex & 懒加载 + fallback 读取。 \\
\bottomrule
\end{longtable}

\section{常见问题排查（FAQ）}
\subsection*{Q1：为什么多选案例后只显示一个？}
你选的是 contour 类图型（FFT / Fluctuation 2D / POD），当前设计仅支持单案例。

\subsection*{Q2：\texttt{vs Gamma\_Linear} 没有曲线？}
请检查：
\begin{itemize}
  \item \texttt{Linear File} 是否有效；
  \item 文件行是否包含可解析的 \(k_y,\omega,\gamma\) 三列；
  \item 终端是否提示找不到文件或无有效行。
\end{itemize}

\subsection*{Q3：POD parity 报并行分辨率不足？}
该模式要求 \(\theta>1\)。若案例本身并行采样过低，程序会阻止绘图并提示 warning。

\subsection*{Q4：rcorr\_phi 没有给出相关长度？}
通常是 SciPy 未安装，或拟合失败。相关曲线仍可显示。

\subsection*{Q5：为什么看不到某些字段（Apar/Bpar）？}
案例可能未输出对应场，或字段维度与当前解析逻辑不匹配。请查看终端中的 shape 提示。

\section{版本维护建议}
\begin{itemize}
  \item 每次新增 Plot Type 后更新本说明书中的“映射表”和“动态选项”章节。
  \item 每次变更默认常量后同步“默认配置”章节。
  \item 建议在项目根目录保存一份示例输入和示例图作为回归检查基线。
\end{itemize}

\section{编译说明}
推荐使用 XeLaTeX：
\begin{lstlisting}
xelatex cgyro_comparison_manual.tex
\end{lstlisting}
如需目录/引用稳定，可连续编译两次。

\end{document}
